\section{Complete fixed-parameter simplification}\label{sec:complexity}

\subsection{Constructively enumerating connected regions}

For completeness, we give the combinatorial ingredients of the incoming
region bound explicitly. This separates the new localization theorem from
the search machinery on which its stronger algorithmic consequence rests.

\begin{lemma}[Connected-region count]\label{lem:regions}
In a graph with $t$ vertices and maximum degree four, the number of connected
$i$-vertex sets is at most $t16^{i-1}$ for $i\ge1$. All connected sets of
size at most $R$ can be enumerated without repetition using polynomial work
per generated set.
\end{lemma}
\begin{proof}
Choose the least vertex as root and a deterministic spanning tree of each
connected set. A depth-first traversal of the tree has $2(i-1)$ edge steps.
Record the root and, at each step, one of at most four local neighbour ports.
The walk reconstructs its visited vertex set, so this is an injective
encoding among the deterministic choices. There are at most
$t4^{2(i-1)}$ such encodings.

To enumerate sets, define the parent of a non-singleton connected set $A$
by deleting its largest non-root vertex whose deletion leaves a connected
set. A spanning tree has a leaf other than the root, so at least one deletion
exists. Iterating the parent reaches the singleton root. From each current
set, consider its at most $4|A|$ frontier neighbours and accept an extension
exactly when its canonical parent is the current set. This visits each set
once. Testing all candidate deletions and their connectivity uses only
polynomial work in $|A|$. An explicit stack gives the implementation without
recursion-depth dependence on the ambient number of vertices.
\end{proof}

The geometric factor $t$ chooses the location of the region. It is not
raised to a power proportional to $U$. This is the step that direct
enumeration of all $R$-subsets would lose.

\subsection{An elementary independent FPT bound}

Before using commutation compression, localization alone already implies
fixed-parameter tractability. In a region containing $i$ original cells,
at most $i+U$ eligible cells are live at any time. A legal eligible event
is specified by a local face or edge of such a cell; there are at most
$10(i+U)$ choices before eliminating repetitions and illegal configurations.
Every first-descent witness has at most $2U+1$ events. Thus direct local
branching, followed over every connected region of size at most $3U+3$,
gives
\begin{equation}\label{eq:elementary-fpt}
 2^{O(U\log(U+2))}\poly(t+B).
\end{equation}
This proof supplies a separate, weaker FPT guarantee that does not require
the constant-alphabet trace theorem. The sharper result below removes the
extra $\log U$ in the exponent.

\subsection{Constant-alphabet encoding of commuting traces}

Assign each eligible incarnation a commitment from
\begin{equation}\label{eq:alphabet}
 \mathcal A=\{\mathrm{idle}\}\ \cup\
 \{\mathrm{down\ at\ edge}\ e:0\le e<6\}\ \cup\
 \{\mathrm{up\ at\ face}\ f:0\le f<4\}.
\end{equation}
An event is ready if all its participants choose their local port for that
event. Two different ready events cannot share an incarnation: a local face
determines its unique paired-face event, and a local edge determines its
unique full edge-star event. Ready events therefore commute by
\cref{lem:forward}.

Fix a deterministic scheduler. If an event is ready, execute the first ready
event; otherwise assign the next unassigned eligible incarnation one of its
eleven symbols. Surviving cells retain their commitments; born cells are
unassigned. Protected original cells are permanently idle.

\begin{lemma}[Commitment encoding]\label{lem:encoding}
For a supplied region with $i$ original cells, every finite trace using
$u$ upward and $d$ downward events has an encoding of length
$i+3u+2d$ over \eqref{eq:alphabet}, up to interchanges of independent events.
Consequently the number of commuting classes with $u\le U$ and
$d\le u+1$ is at most
\begin{equation}\label{eq:descent-count}
 11^{i+5U+2}.
\end{equation}
The generic bounded-region family, without the latter loss restriction,
has at most $11^{3i+5U}$ classes.
\end{lemma}
\begin{proof}
For a target trace, give each incarnation the port of the event that
eventually consumes it, and give every target survivor the symbol idle.
Follow these choices when the deterministic scheduler asks for them.

If an event is ready, consider the first remaining target event touching
one of its source incarnations. All earlier events leave those sources
intact. Its commitment identifies exactly that ready event, so the first
touch must be the ready event itself. All preceding target events have
disjoint consumed supports; the persistence and relative-boundary lemmas
move the ready event to the front. Iterating realizes the target trace
up to commuting independent events. Assigning the recorded symbols cannot
stop before the next target event is ready, and target survivors are idle.

The freely assigned incarnations comprise the $i$ eligible originals and
the $3u+2d$ born cells. Each receives one symbol. Record assignments in the
order requested by the deterministic scheduler. Padding by idle symbols
to a common maximum length preserves an injective encoding of trace
classes: decoding a word determines one execution. Substituting
$d\le u+1$ yields \eqref{eq:descent-count}. In a general region-supported
trace, protected originals survive, so $d\le i+u$, giving the final bound.
\end{proof}

The counting decoder in this lemma is allowed to pass an earlier descent
introduced by reordering. It need not follow the implementation's early-stop
policy. The target trace's event counts and total encoding length are what
establish the injection. Confusing these two roles would unnecessarily weaken
or invalidate the counting argument.

The practical search uses sleep sets instead of literal assignments. It
explores enabled events in deterministic order and remembers an earlier
sibling event across a later branch only while its consumed incarnation
set is disjoint from each intervening event. Once a source incarnation is
destroyed it never reappears. This nonresurrection, local event uniqueness,
and persistence imply that at most one representative of a commuting class
is visited. Conversely, the first representative of each class in the fixed
search order cannot be blocked: a blocking event could be commuted left to
obtain an earlier representative. This is the exact finite sleep-set argument
specialized to these geometric events, as implemented in the incoming engine
\cite{incoming-trace}.

\begin{theorem}[Complete bounded-upward descent]\label{thm:fpt}
For a valid generalized source $T$ of size $t\ge1$, an optional integral
marking of input bit bound $B$, and $U\ge0$, question \eqref{eq:decision}
is decidable in
\begin{equation}\label{eq:fpt}
 t\,2^{O(U)}\poly(t+U+B)=2^{O(U)}\poly(t+B)
\end{equation}
bit operations. A positive answer supplies a geometrically replayable
first-descent trace. No spatial-coverage hypothesis is required.
\end{theorem}
\begin{proof}
Set $R=\min(t,3U+3)$. Enumerate all connected original regions of size at
most $R$ and search their bounded-upward traces with the endpoint predicate
$|T'|<t$. A positive test returns immediately. If no positive is found,
every permitted family is exhausted. The localization theorem proves that
this decides \eqref{eq:decision} completely.

Every visited prefix before success has $d\le u$, and the first successful
one has $d=u+1$. Combining \cref{lem:regions,lem:encoding}, a sufficient
counting factor, including the source endpoint, is
\begin{equation}\label{eq:explicit-count}
 1+t\sum_{i=1}^{R}16^{i-1}11^{i+5U+2}.
\end{equation}
Since $R\le3U+3$, this is $t2^{O(U)}$ times a fixed constant. In particular,
the largest commitment exponent is $8U+5$, improving the $14U+9$ obtained
by substituting $R$ into the generic region bound. The bit analysis below
supplies polynomial work per visited encoding or representative. It includes
the endpoint count test. This proves \eqref{eq:fpt}.
\end{proof}

\begin{corollary}[Small-budget regimes]\label{cor:qp}
Bounded-upward descent is polynomial-time when $U=O(\log(t+2))$, and takes
$t^{O(\log(t+2))}\poly(B)$ time when $U=O(\log^2(t+2))$.
\end{corollary}
\begin{proof}
Substitute the budget into \eqref{eq:fpt}. Polynomial factors in $U$ are
absorbed by the indicated bounds. The bit bound $B$ remains polynomially
charged.
\end{proof}

The constants in \eqref{eq:explicit-count} are intentionally conservative.
It is an asymptotic guarantee, not a prediction that large $U$ is practical.
For $U=0$, the problem is simply the existence of a legal initial downward
move, since any descending trace has a first downward move. Directly checking
edge stars solves that special case without a general region enumeration.

\subsection{Binary arithmetic and certificate size}

\begin{lemma}[Upward moves control marking growth]\label{lem:bits}
Let $D_0$ be the maximum absolute source edge value. After a trace with $u$
upward events, all signed edge values have magnitude at most $2^uD_0$.
Normalized height rows and local coherent normal coordinates have bit
length at most $B+u+O(1)$.
\end{lemma}
\begin{proof}
A downward move removes one edge and creates no new edge; all surviving
edge values remain unchanged. The new diagonal of an upward move joins
the two formal apices $a,b$. Through any belt vertex $v$, its value is
$(h_b-h_v)+(h_v-h_a)$, a sum of two old signed edge values. Its magnitude
is at most twice the previous maximum, and all other edge values persist.
Induct on $u$. A normalized row entry is an edge difference from its first
corner. A local normal coordinate is a gap between ordered corner heights,
bounded by their span, which is itself an edge difference.
\end{proof}

There are $O(t+U)$ local edge and face occurrences at every visited state.
Global manifold validation, grouping edge stars, comparing local patterns,
and copying a full triangulation therefore cost polynomial time in
$t+U+B$. The code currently rebuilds global geometry where needed; it makes
no constant-time local-update claim. Stable incarnation names are exact
interned DAG expressions, not cryptographic identity assumptions.

A first-descent proof with complete intermediate snapshots has bit length
\begin{equation}\label{eq:certificate-size}
 O\bigl((U+1)(t+U)(B+U+\log(t+U+1))\bigr).
\end{equation}
It contains at most $2U+1$ moves, each with at most $t+U$ rows, plus the
final coordinate data. Independent replay is polynomial in this size.
Exhaustive failure is an algorithmic conclusion about a bounded family;
the implementation does not claim a comparably short negative certificate.

\subsection{Repeated bounded simplification}

\begin{definition}
A triangulation is $U$-irreducible for this move model if it admits no
strict descent using at most $U$ total upward moves.
\end{definition}

\begin{corollary}[Complete bounded preprocessing]\label{cor:preprocess}
Repeatedly applying complete bounded-$U$ descent and adopting each verified
positive result terminates at a $U$-irreducible triangulation in
$2^{O(U)}\poly(t+B)$ bit time. Its concatenated positive proof uses at most
$tU$ upward events and $t(2U+1)$ events in total.
\end{corollary}
\begin{proof}
Each successful restart reduces the tetrahedron count by at least one,
so there are fewer than $t$ successful restarts. Each local transient state
has at most $t+U$ tetrahedra. Across all successful traces, \cref{lem:bits}
increases the marking bit bound by at most $tU+O(1)$. Summing the polynomial
overheads in \cref{thm:fpt} gives the stated bound. The final complete
negative bounded search is precisely the definition of $U$-irreducibility.
\end{proof}

An arbitrary choice of successful one-unit descents is not guaranteed to
preserve a prescribed total upward budget for a later $k$-unit objective.
The next section treats that joint-budget problem separately.
